{{دیوان
| عنوان = پرندے کی فریاد
| ماخذ = ویکی ماخذ
| ماخذ_لنک = https://ur.wikisource.org/wiki/%D9%BE%D8%B1_%D9%86%D8%AF%DB%92_%DA%A9%DB%8C_%D9%81%D8%B1%DB%8C%D8%A7%D8%AF
| مقابلہ = اقبال سائبر لائبریری (اقبال اکادمی پاکستان)
| مقابلہ_لنک = https://www.iqbalcyberlibrary.net
}}
=== بچوں کے لیے ===

<poem>
آتا ہے یاد مجھ کو گزرا ہوا زمانا
وہ باغ کی بہاریں، وہ سب کا چہچہانا

آزادیاں کہاں وہ اب اپنے گھونسلے کی
اپنی خوشی سے آنا، اپنی خوشی سے جانا

لگتی ہے چوٹ دل پر، آتا ہے یاد جس دم
شبنم کے آنسوؤں پر کلیوں کا مُسکرانا

وہ پیاری پیاری صورت، وہ کامنی سی مورت
آباد جس کے دم سے تھا میرا آشیانا

آتی نہیں صدائیں اُس کی مرے قفس میں
ہوتی مری رہائی اے کاش میرے بس میں!

کیا بد نصیب ہوں مَیں گھر کو ترس رہا ہوں
ساتھی تو ہیں وطن میں، مَیں قید میں پڑا ہوں

آئی بہار، کلیاں پھُولوں کی ہنس رہی ہیں
مَیں اس اندھیرے گھر میں قسمت کو رو رہا ہوں

اس قید کا الٰہی! دُکھڑا کسے سُناؤں
ڈر ہے یہیں قفسں میں مَیں غم سے مر نہ جاؤں

جب سے چمن چھُٹا ہے، یہ حال ہو گیا ہے
دل غم کو کھا رہا ہے، غم دل کو کھا رہا ہے

گانا اسے سمجھ کر خوش ہوں نہ سُننے والے
دُکھّے ہوئے دلوں کی فریاد یہ صدا ہے

آزاد مجھ کو کر دے، او قید کرنے والے!
مَیں بے زباں ہوں قیدی، تو چھوڑ کر دُعا لے
</poem>
